% =============================================================
%  Chapter -- DC Line
%  Subsections: Resistance, Voltage drop and losses, Current for a
%  given power, Worked example.
%  Values: link/fig/link_values.tex (generated by link/fig/link_example.py
%  from tools/example_data.py).
%  Sources: [1] hvdc_lf, src/hvdc_lf/components/dc_line/calc.py.
% =============================================================
\section{DC Line}
\label{sec:dc-line}

The DC line connects the DC terminals of the two stations of a link. One pole conductor is calculated; a pole may
consist of $n_{sub}$ conductors in parallel. The line is modelled by its DC resistance; inductance and capacitance
have no effect in the steady state. Both pole conductors carry the same current $I_d$ (balanced operation, no current
in the return path), so the line losses of the link are $n_{pc}$ times those of one pole conductor, $n_{pc} = 2$.

\subsection{Resistance}

The conductor data give the DC resistance per kilometre $R'_{ref}$ at the reference temperature $T_{ref}$ and its
temperature coefficient $\alpha$. The resistance of one pole at the conductor temperature $T$ is
\begin{equation}
  R_{pole}(T) = \frac{R'_{ref} \cdot \left( 1 + \alpha \cdot (T - T_{ref}) \right)}{n_{sub}} \cdot l .
\end{equation}
The conductor temperature is an operating value; without one the design temperature $T_{design}$ of the line is
used.

\subsection{Voltage Drop and Losses}

With the voltage $U_{send}$ pole-to-ground at the sending end (rectifier DC terminal) and the pole current $I_d$,
\begin{equation}
  U_{recv} = U_{send} - R_{pole} \cdot I_d , \qquad
  P_{send} = U_{send} \cdot I_d , \qquad
  P_{recv} = U_{recv} \cdot I_d , \qquad
  P_{loss\_line} = R_{pole} \cdot I_d^2 .
\end{equation}
These powers are per pole and taken as positive in the direction of the power flow. The loading of the line is
$I_d / (n_{sub} \cdot I_r)$ with the rated current $I_r$ of one conductor, when it is given.

\subsection{Current for a Given Power}

For a power given at the sending end, $I_d = P_{send} / U_{send}$. For a power given at the receiving end,
$P_{recv} = (U_{send} - R_{pole} \cdot I_d) \cdot I_d$ is a quadratic equation in $I_d$ with the smaller root
\begin{equation}
  I_d = \frac{U_{send} - \sqrt{U_{send}^2 - 4 \cdot R_{pole} \cdot P_{recv}}}{2 \cdot R_{pole}} .
\end{equation}
A real solution exists up to $P_{recv\_max} = U_{send}^2 / (4 \cdot R_{pole})$; a larger power is rejected. The
larger root belongs to operation beyond this maximum, with $U_{recv} < U_{send} / 2$, and is not used.

\subsection{Worked Example}
\label{sec:dc-line-example}

The example line is a 300\,km XLPE cable, 2500\,mm$^2$ copper, one conductor per pole, $R'_{ref} = 0.0072\,\Omega$/km
at 20\,$^{\circ}$C, $\alpha = 0.00393$\,1/K, $I_r = 2.0$\,kA and $T_{design} = 70\,^{\circ}$C. The pole resistance is
$\lkRtwenty\,\Omega$ at 20\,$^{\circ}$C and $\lkRseventy\,\Omega$ at 70\,$^{\circ}$C. At $U_{send} = 400$\,kV and
$I_d = \lkId$\,kA the receiving end voltage is $\lkUinv$\,kV and the losses of both pole conductors are
$\lkLine$\,MW. The maximum power at the receiving end is 15\,480\,MW per pole, about 31 times the power of the example.
